\problem{3.41}{Largest momentum expectation for two oscillator levels}
Equal probabilities for the first two energies imply
the general normalized state
\begin{equation}
\ket{\Psi(0)}=\frac1{\sqrt2}
(\ket0+\ee^{\ii\theta}\ket1),
\end{equation}
where a global phase has been removed.
Because $p=\ii\sqrt{m\hbar\omega/2}(a^\dagger-a)$
and $\avg a=\ee^{\ii\theta}/2$,
\begin{equation}
\boxed{\avg p(0)=
\sqrt{\frac{m\hbar\omega}{2}}\sin\theta}.
\end{equation}
The largest \emph{positive} value is
$\boxed{\sqrt{m\hbar\omega/2}}$,
achieved at $\theta=\pi/2$.
Including time dependence,
\begin{equation}
\boxed{\Psi(x,t)=\frac1{\sqrt2}
\left[\psi_0(x)\ee^{-\ii\omega t/2}+
\ii\psi_1(x)\ee^{-3\ii\omega t/2}\right]}.
\end{equation}
The momentum expectation later evolves
as $\avg p(t)=\sqrt{m\hbar\omega/2}\cos\omega t$.

\problem{3.42}{Coherent states, their normalization and time evolution}
Let $a\ket\alpha=\alpha\ket\alpha$,
where $\alpha$ is arbitrary complex.
Taking the adjoint gives
$\bra\alpha a^\dagger=\alpha^*\bra\alpha$.
For $x=\ell_x(a+a^\dagger)$,
$p=\ii\ell_p(a^\dagger-a)$,
with $\ell_x=\sqrt{\hbar/(2m\omega)}$ and
$\ell_p=\sqrt{m\hbar\omega/2}$,
\begin{align}
\boxed{\avg x}&=\boxed{\ell_x(\alpha+\alpha^*)},&
\boxed{\avg p}&=\boxed{\ii\ell_p(\alpha^*-\alpha)},\\
\boxed{\avg{x^2}}&=\boxed{\avg x^2+\ell_x^2},&
\boxed{\avg{p^2}}&=\boxed{\avg p^2+\ell_p^2}.
\end{align}
Indeed $\avg{a^\dagger a}=|\alpha|^2$,
$\avg{aa^\dagger}=|\alpha|^2+1$,
$\avg{a^2}=\alpha^2$ and
$\avg{(a^\dagger)^2}=(\alpha^*)^2$.
It follows that
\begin{equation}
\boxed{\sigma_x=\sqrt{\frac{\hbar}{2m\omega}}},
\quad
\boxed{\sigma_p=\sqrt{\frac{m\hbar\omega}{2}}},
\quad
\boxed{\sigma_x\sigma_p=\frac\hbar2}.
\end{equation}
Expand $\ket\alpha=\sum_{n=0}^\infty c_n\ket n$.
Comparing coefficients in
$a\ket\alpha=\alpha\ket\alpha$ yields
$c_{n+1}\sqrt{n+1}=\alpha c_n$, hence
\begin{equation}
c_n=\frac{\alpha^n}{\sqrt{n!}}c_0.
\end{equation}
Normalization gives
\begin{equation}
1=|c_0|^2\sum_{n=0}^\infty
\frac{|\alpha|^{2n}}{n!}
=|c_0|^2\ee^{|\alpha|^2}
\quad\Rightarrow\quad
\boxed{c_0=\ee^{-|\alpha|^2/2}}
\end{equation}
up to an overall phase.
Under oscillator evolution,
\begin{equation}
\ket{\Psi(t)}=
\ee^{-\ii\omega t/2}
\ket{\alpha\ee^{-\ii\omega t}},
\quad
\boxed{\alpha(t)=\alpha\ee^{-\ii\omega t}}.
\end{equation}
Thus the state remains coherent, with fixed
uncertainties. If
$\alpha=C\sqrt{m\omega/(2\hbar)}
\ee^{\ii\phi}$, then
\begin{equation}
\boxed{\avg{x}(t)=C\cos(\omega t-\phi)},
\qquad
\boxed{\sigma_x(t)=\sqrt{\frac{\hbar}{2m\omega}}}.
\end{equation}
The ground state $\ket0$ is coherent with
$\boxed{\alpha=0}$.

\problem{3.43}{Robertson–Schrödinger uncertainty bound}
For Hermitian $A,B$, set
$\delta A=A-\avg A\,I$,
$\delta B=B-\avg B\,I$.
Cauchy--Schwarz gives
\begin{equation}
\sigma_A^2\sigma_B^2\ge
|\avg{\delta A\delta B}|^2.
\end{equation}
Split the product into symmetric
and antisymmetric pieces:
\begin{equation}
\avg{\delta A\delta B}=
\frac12\avg{\{\delta A,\delta B\}}+
\frac12\avg{[A,B]}.
\end{equation}
Let $C=-\ii[A,B]$ and
$D=AB+BA-2\avg A\avg B\,I$.
Then $C,D$ are Hermitian, and
$\avg{[A,B]}=\ii\avg C$,
$\avg{\{\delta A,\delta B\}}=\avg D$.
Thus
\begin{equation}
\boxed{\sigma_A^2\sigma_B^2\ge
\frac14\left(\avg C^2+\avg D^2\right)},
\end{equation}
where $\avg C^2$ means
$(\avg C)^2$, not $\avg{C^2}$.
For $B=A$, $C=0$ and
$\avg D=2(\avg{A^2}-\avg A^2)=2\sigma_A^2$.
The extended inequality becomes
$\sigma_A^4\ge\sigma_A^4$:
it is automatically an equality and
adds no new information.

\problem{3.44}{Time evolution of a coupled three-level Hamiltonian}
In the basis $(e_1,e_2,e_3)$,
\begin{equation}
H=\begin{pmatrix}a&0&b\\0&c&0\\b&0&a\end{pmatrix}.
\end{equation}
The vector $e_2$ is an eigenstate
with energy $c$. For initial state $e_2$,
\begin{equation}
\boxed{\ket{S(t)}=\ee^{-\ii ct/\hbar}e_2}.
\end{equation}
The remaining block is
$aI+b\sigma_x$ on $\operatorname{span}(e_1,e_3)$.
Its eigenvectors $(e_1\pm e_3)/\sqrt2$
have energies $a\pm b$.
Exponentiation in that subspace gives
\begin{align}
\ee^{-\ii t(aI+b\sigma_x)/\hbar}
&=\ee^{-\ii at/\hbar}
[\cos(bt/\hbar)I-\ii\sin(bt/\hbar)\sigma_x].
\end{align}
Thus for initial state $e_1$,
\begin{equation}
\boxed{\ket{S(t)}=\ee^{-\ii at/\hbar}
\left[\cos\frac{bt}{\hbar}\,e_1
-\ii\sin\frac{bt}{\hbar}\,e_3\right]}.
\end{equation}
Its two nonzero probabilities are
$\cos^2(bt/\hbar)$ and $\sin^2(bt/\hbar)$,
summing to one.

\problem{3.45}{Position operator acting on energy coefficients}
If $\ket{S(t)}=\sum_{k=0}^\infty c_k(t)\ket k$,
insert the oscillator position matrix elements
\begin{equation}
\bra n x\ket k=
\sqrt{\frac{\hbar}{2m\omega}}
[\sqrt{k}\,\delta_{n,k-1}
+\sqrt{k+1}\,\delta_{n,k+1}].
\end{equation}
Summing over $k$ gives
\begin{equation}
\boxed{\bra n x\ket{S(t)}=
\sqrt{\frac{\hbar}{2m\omega}}
\left[\sqrt n\,c_{n-1}(t)
+\sqrt{n+1}\,c_{n+1}(t)\right]}.
\end{equation}
Set $c_{-1}=0$. The position operator
therefore mixes only neighboring energy levels;
it is not diagonal in the energy basis.

\problem{3.46}{Three-level spectra and outcome probabilities}
Let $e_1,e_2,e_3$ be the standard basis.
The specified observables are
\begin{equation}
H=\hbar\omega\,\mathrm{diag}(1,2,2),\quad
A=\lambda\begin{pmatrix}0&1&0\\1&0&0\\0&0&2\end{pmatrix},\quad
B=\mu\begin{pmatrix}2&0&0\\0&0&1\\0&1&0\end{pmatrix}.
\end{equation}
An orthonormal eigenbasis for $H$ is
$\{e_1,e_2,e_3\}$, energies
$\hbar\omega,2\hbar\omega,2\hbar\omega$.
For $A$, eigenvalues and eigenstates are
\begin{equation}
\begin{array}{c|c}
+\lambda&(e_1+e_2)/\sqrt2\\
-\lambda&(e_1-e_2)/\sqrt2\\
2\lambda&e_3
\end{array}.
\end{equation}
For $B$ they are
\begin{equation}
\begin{array}{c|c}
2\mu&e_1\\
+\mu&(e_2+e_3)/\sqrt2\\
-\mu&(e_2-e_3)/\sqrt2
\end{array}.
\end{equation}
For $\ket{S(0)}=c_1e_1+c_2e_2+c_3e_3$,
$\sum|c_j|^2=1$, the expectations are
\begin{align}
\boxed{\avg H}&=\hbar\omega
(|c_1|^2+2|c_2|^2+2|c_3|^2),\\
\boxed{\avg A}&=2\lambda
[\operatorname{Re}(c_1^*c_2)+|c_3|^2],\\
\boxed{\avg B}&=2\mu
[|c_1|^2+\operatorname{Re}(c_2^*c_3)].
\end{align}
Time evolution is
\begin{equation}
\boxed{\ket{S(t)}=
c_1\ee^{-\ii\omega t}e_1+
(c_2e_2+c_3e_3)\ee^{-2\ii\omega t}}.
\end{equation}
The energy probabilities remain
$P(\hbar\omega)=|c_1|^2$ and
$P(2\hbar\omega)=|c_2|^2+|c_3|^2$.
For $A$ the probabilities are
\begin{align}
P(A=+\lambda)&=\tfrac12|c_1\ee^{\ii\omega t}+c_2|^2,\\
P(A=-\lambda)&=\tfrac12|c_1\ee^{\ii\omega t}-c_2|^2,\\
P(A=2\lambda)&=|c_3|^2.
\end{align}
For $B$, the common phase in the
degenerate energy subspace cancels:
\begin{align}
P(B=2\mu)&=|c_1|^2,\\
P(B=+\mu)&=\tfrac12|c_2+c_3|^2,\\
P(B=-\mu)&=\tfrac12|c_2-c_3|^2.
\end{align}
Each distribution sums to one.

\problem{3.47}{Supersymmetric partner Hamiltonians and delta wells}
Write $s=\hbar/\sqrt{2m}$.
Since $p=-\ii\hbar\partial_x$,
\begin{equation}
A=s\partial_x+W(x),\quad
A^\dagger=-s\partial_x+W(x).
\end{equation}
For real $W$ and suitable domains,
multiplying and applying the product rule gives
\begin{align}
H_1=A^\dagger A&=-s^2\partial_x^2+
W^2-sW',\\
H_2=AA^\dagger&=-s^2\partial_x^2+
W^2+sW'.
\end{align}
Thus
\begin{equation}
\boxed{V_1=W^2-sW',\qquad V_2=W^2+sW'}.
\end{equation}
If $H_1\psi_n^{(1)}=E_n\psi_n^{(1)}$,
then
$H_2(A\psi_n^{(1)})=A H_1\psi_n^{(1)}
=E_n(A\psi_n^{(1)})$.
Likewise $H_1(A^\dagger\psi_n^{(2)})
=E_n(A^\dagger\psi_n^{(2)})$.
A nonzero energy eigenstate maps to a
partner state unless annihilated by the
intertwiner. If $A\psi_0^{(1)}=0$, then
\begin{equation}
\boxed{W(x)=-s\,\frac{\psi_0'(x)}{\psi_0(x)}}.
\end{equation}
For the shifted attractive delta well,
$\psi_0(x)=\sqrt\kappa\,\ee^{-\kappa|x|}$,
$\kappa=m\alpha/\hbar^2$, hence
\begin{equation}
W(x)=s\kappa\,\operatorname{sgn}(x)
=\alpha\sqrt{\frac{m}{2\hbar^2}}
\,\operatorname{sgn}(x).
\end{equation}
Use $(\operatorname{sgn}x)'=2\delta(x)$:
\begin{equation}
W^2=\frac{m\alpha^2}{2\hbar^2},\qquad
sW'=\alpha\delta(x).
\end{equation}
Therefore
\begin{equation}
\boxed{V_1=\frac{m\alpha^2}{2\hbar^2}
-\alpha\delta(x),\qquad
V_2=\frac{m\alpha^2}{2\hbar^2}
+\alpha\delta(x)}.
\end{equation}
The partner is a \emph{repulsive} delta
plus the same constant shift, so
it has no negative-energy bound level.
Its continuum scattering probabilities
equal those of the attractive partner.

\problem{3.48}{Symmetry versus self-adjointness of momentum}
Consider $p=-\ii\hbar\partial_x$ on $(0,a)$.
For differentiable $f,g$,
integration by parts gives
\begin{equation}
\braket g{pf}-\braket{pg}f
=-\ii\hbar[g^*(x)f(x)]_0^a.
\end{equation}
(a) For Dirichlet boundary conditions
$f(0)=f(a)=0$ this vanishes for every $g$
without requiring $g$ to vanish at the
endpoints. The operator is symmetric on
$H_0^1(0,a)$, but its adjoint domain is
larger (essentially $H^1(0,a)$), so it is
not self-adjoint.
(b) If $f(a)=\lambda f(0)$,
the boundary term vanishes against all
domain vectors $f$ precisely when
$\lambda g^*(a)-g^*(0)=0$,
or
\begin{equation}
\boxed{g(a)=\frac{g(0)}{\lambda^*}}
\quad(\lambda\ne0).
\end{equation}
For the adjoint domain to be identical
to the operator domain, require
$\lambda=1/\lambda^*$, equivalently
$\boxed{|\lambda|=1}$.
Thus all phases $\lambda=\ee^{\ii\theta}$
give self-adjoint quasi-periodic momentum
operators on the finite interval.
(c) On $(0,\infty)$, the deficiency equations
$p^\dagger g=\pm\ii\gamma g$, $\gamma>0$,
give
$g_+(x)\propto\ee^{-\gamma x/\hbar}$,
$g_-(x)\propto\ee^{+\gamma x/\hbar}$.
Only the first is square integrable,
so the deficiency indices are unequal:
$\boxed{(n_+,n_-)=(1,0)}$ for this convention.
Therefore there is no self-adjoint
extension of $p$ on the half-line.
\qnote{Momentum domains in a finite box
are not determined solely by the
formal differential expression.}

\problem{3.49}{Gaussian free-particle evolution in momentum space}
(a) For $H=p^2/(2m)$ the momentum-space
Schrödinger equation is an ODE at each $p$:
\begin{equation}
\ii\hbar\partial_t\Phi(p,t)=
\frac{p^2}{2m}\Phi(p,t)
\quad\Longrightarrow\quad
\boxed{\Phi(p,t)=
\Phi(p,0)\ee^{-\ii p^2t/(2m\hbar)}}.
\end{equation}
(b) For the traveling Gaussian
$\Psi(x,0)=(2a/\pi)^{1/4}
\ee^{-ax^2+\ii\ell x}$ with $a>0$,
completing the Fourier-transform square yields
\begin{equation}
\boxed{\Phi(p,0)=
(2\pi a\hbar^2)^{-1/4}
\exp\left[-\frac{(p-\hbar\ell)^2}
{4a\hbar^2}\right]}.
\end{equation}
Therefore $|\Phi(p,t)|^2$ is the
time-independent Gaussian
\begin{equation}
\boxed{|\Phi(p,t)|^2=
\frac1{\sqrt{2\pi a\hbar^2}}
\exp\left[-\frac{(p-\hbar\ell)^2}
{2a\hbar^2}\right]}.
\end{equation}
(c) Standard Gaussian integration gives
\begin{equation}
\boxed{\avg p=\hbar\ell,\quad
\avg{p^2}=\hbar^2(\ell^2+a),\quad
\sigma_p=\hbar\sqrt a}.
\end{equation}
(d) The energy expectation decomposes as
\begin{equation}
\boxed{\avg H=\frac{\hbar^2(\ell^2+a)}{2m}
=\frac{\avg p^2}{2m}
+\underbrace{\frac{\hbar^2 a}{2m}}_{\avg H_0}},
\end{equation}
where $\avg p^2$ in the first fraction
of the final expression denotes
$(\avg p)^2$, not $\avg{p^2}$.
The first term is center-of-mass
translation energy and the second is
the internal spreading energy of
the Gaussian wavepacket.
